\documentclass{article}
\usepackage{/home/user1/code/tex/sty}
\usepackage{times}
\usepackage{amsfonts}
\usepackage{amsmath}
\usepackage{geometry}
\usepackage{pdfpages}
\usepackage{hyperref}
\setlength{\parindent}{0pt}
\geometry{top=1in,bottom=1in,left=1in,right=1in}

\usepackage{graphicx}
\usepackage{floatrow}
\newcommand{\fig}[3]{
\begin{figure}[H]
\begin{center}
\includegraphics[height=#1]{#2}
\caption{#3}
\label{#2}
\end{center}
\end{figure}}
%\fig{5cm}{1.png}{Diagram.}

\begin{document}
\begin{center}
\textbf{\Large{Problem Set 4}}\\~\\
\begin{tabular}{l l}
Student&Agustin Romero\\
Email&ar1@stanford.edu\\
Instructor&Roger Blandford\\
Assistant&Andrew Sullivan\\
Course&PHYSICS 262\\
Institution&Stanford University\\
Date&\today\\
\end{tabular}
\end{center}
\section*{Problem 1}
\section*{Problem 2.a.}
Starting with the definition
\g{p^\alpha = a A^\alpha + b B^\alpha}
apply two covariant derivatives and simplify.
\n{aal}{A^\alpha}
\n{bal}{B^\alpha}
\g{\zeich{ns8na}p^\alpha_{;\gamma \delta}&=(a\aal +b\bal)_{;\gamma \delta}=(a\aal)_{;\gamma \delta}+(b\bal)_{;\gamma \delta}\\
&= (a_{;\gamma}\aal + a\aal_{;\gamma})_{;\delta} + (b_{;\gamma}\bal + a\bal_{;\gamma})_{;\delta}}
This is identical to $p^{\alpha}_{;\delta \gamma}$ except the reversal of $\delta$ and $\gamma$.
\g{\zeich{nz8c}p^{\alpha}_{;\delta \gamma}=(a_{;\delta}\aal + a\aal_{;\delta})_{;\gamma} + (b_{;\delta}\bal + a\bal_{;\delta})_{;\gamma}}
Now the difference between \gz{ns8na} and \gz{nz8c} is
\g{p^\alpha_{;\gamma \delta}-p^{\alpha}_{;\delta \gamma}&=\aal(a_{;\gamma \delta}-a_{;\delta \gamma})+\bal(b_{;\gamma \delta}-b_{;\delta \gamma})+a(\aal_{;\gamma \delta}-\aal_{\delta \gamma})+b(\bal_{\gamma \delta}-B_{\delta \gamma})
\intertext{Covariant derivatives on scalars commute, $a_{;\gamma \delta}=a_{;\delta \gamma}$. The two terms with scalar derivatives cancel.}
&= \boxed{a(\aal_{;\gamma \delta}-\aal_{\delta \gamma})+b(\bal_{\gamma \delta}-B_{\delta \gamma})}}
Giving the desired result.
\section*{Problem 2.b.}
c.f. Lecture 10. For a general vector field $V$
\g{\zeich{nzx8cn}V^\alpha_{;\gamma \delta}-V^\alpha_{;\delta \gamma}=-{R^\alpha}_{\beta \gamma \delta}V^\beta}
So the goal is to calculate $V^\alpha_{;\gamma \delta}$ and simplify. As requested we will start with the first directional derivative of $p$
\g{{p^\alpha}_{;\gamma}={p^\alpha}_{,\gamma}+{p^\beta}\Gamma^\alpha_{\beta \gamma}}
We will also need the covariant derivative of $\Gamma^{\alpha}_{\mu\gamma}$
\g{\Gamma^\alpha_{\mu\gamma;\delta}=\Gamma^\alpha_{\mu\gamma,\delta}+\Gamma^\alpha_{\nu\delta}\Gamma^\nu_{\mu\delta}-\Gamma^\nu_{\mu\delta}\Gamma^\alpha_{\nu\delta}-\Gamma^\nu_{\delta\gamma}\Gamma^\alpha_{\mu\nu}}
Take the second directional derivative of ${p^\alpha}_{;\gamma}$,
\g{p^\alpha_{;\gamma\delta}&=(p^\alpha_{,\gamma}+\Gamma^\alpha_{\mu\gamma}p^\mu)_{;\delta}\\
&=p^\alpha_{,\gamma\delta}+\Gamma^\alpha_{\mu\delta}p^\mu_{,\gamma}-\Gamma^\mu_{\gamma\delta}p^\alpha_{,\mu}+\Gamma^{\alpha}_{\mu\gamma}p^\mu_{,\gamma}+\Gamma^\alpha_{\mu\gamma}\Gamma^\mu_{\nu\delta}p^\nu+\Gamma^\alpha_{\mu\gamma;\delta}p^\mu
\intertext{Notice some terms have $p^\nu$ and $p^\mu$. They will be relabled later. Expand the double derivative on $\Gamma$, the last term. All other terms remain the same.}
&=p^\alpha_{,\gamma\delta}+\Gamma^\alpha_{\mu\delta}p^\mu_{,\gamma}-\Gamma^\mu_{\gamma\delta}p^\alpha_{,\mu}+\Gamma^{\alpha}_{\mu\gamma}p^\mu_{,\gamma}+\Gamma^\alpha_{\mu\gamma}\Gamma^\mu_{\nu\delta}p^\nu+(\Gamma^\alpha_{\mu\gamma,\delta}+\Gamma^\alpha_{\nu\delta}\Gamma^\nu_{\mu\delta}-\Gamma^\nu_{\mu\delta}\Gamma^\alpha_{\nu\delta}-\Gamma^\nu_{\delta\gamma}\Gamma^\alpha_{\mu\nu})p^\mu}
This is our key equation. We know that we only need to swap $\gamma$ and $\delta$ to then calculate \gz{nzx8cn}. We can see what will transpire: 1. The term which is a double partial derivative $p^\alpha_{,\gamma\delta}$ will cancel by commutivity, 2) The term which contains $\Gamma^\mu_{\gamma\delta}$ will cancel due to the commutivity of the lower indices. Now expand,
\g{p^\alpha_{;\gamma\delta}-p^\alpha_{;\delta\gamma}&=(\Gamma^\alpha_{\mu\gamma}\Gamma^\mu_{\nu\delta}-\Gamma^\alpha_{\mu\delta}\Gamma^\mu_{\nu\delta})p^\nu + (\Gamma^\alpha_{\mu\gamma,\delta}-\Gamma^\alpha_{\mu\delta,\gamma}+\Gamma^\alpha_{\nu\delta}\Gamma^\nu_{\mu\gamma}-\Gamma^\alpha_{\nu\delta}\Gamma^\nu_{\mu\delta}-\Gamma^\nu_{\mu\delta}\Gamma^\alpha_{\nu\gamma}+\Gamma^\nu_{\mu\delta}\Gamma^\alpha_{\nu\delta})p^\mu
\intertext{This looks daunting, however note that the $p^\nu$ term can be relabled with $\nu\rightarrow \mu$. This immediately cancels the middle two terms in the $p^\mu$ parenthesis. The final result is}
&= (\Gamma^\alpha_{\mu\gamma,\delta}-\Gamma^\alpha_{\mu\delta,\gamma}-\Gamma^\nu_{\mu\delta}\Gamma^\alpha_{\nu\gamma}+\Gamma^\nu_{\mu\delta}\Gamma^\alpha_{\nu\delta})p^\mu}
This is indeed the desired form of $R^\alpha_{\beta\gamma\delta}$ but the arrangement is unclear. Relablel $\nu\rightarrow\epsilon$ and $\delta\rightarrow\gamma$ again, to obtain the form
\g{\boxed{R^\alpha_{\beta\gamma\delta}=\Gamma^\alpha_{\beta\delta,\gamma}-\Gamma^\alpha_{\beta\gamma,\delta}-\Gamma^\epsilon_{\beta\gamma}\Gamma^\alpha_{\epsilon\gamma}+\Gamma^\epsilon_{\beta\delta}\Gamma^\alpha_{\epsilon\delta}}}
Which is the desired form of $R$.
\section*{Problem 3.a.}
The solution for question 3 is provided in the following Mathematic pages.
\section*{Problem 3.b.}
Included in the following Mathematica pages. As a comment, the Einstein tensor at the throat is dependent on $\theta$ and symmetric about $\theta=0$.
\section*{Problem 3.c.}
Included in the following Mathematica pages.
\section*{Problem 3.d.}
Included in the following Mathematica pages.
\includepdf[pages=-]{q3.pdf}
\section*{Problem 4.a.}
\includepdf[pages=-]{q4.pdf}
\section*{Problem 4.b.}
\section*{Problem 4.c.}
\section*{Problem 4.d.}
\section*{Problem 4.d.}
\end{document}
